% DO NOT UPLOAD. First author-style draft, 2026-09-07.
% Classification: PAPER-CANDIDATE. Not an accepted resolution or release.
% Sole mathematical source: ../certificate_r1/, manifest 4076 bytes, SHA256
% 89992CA4E10051B3C0E953B7E09E11EABD643459CBF2F86B7173D745D089EBCA.
% Manuscript equation cross-check: PENDING independent review.
% Style: supplied Haraux/ABPL/IBPL/CJE sources; no venue or authorship chosen.
% Source IDs D1-D10, C1-C12, L1-L7 are retained as equation tags.
\documentclass[11pt]{article}
\usepackage{amsmath,amssymb,amsthm}
\setlength{\textwidth}{6.3in}
\setlength{\oddsidemargin}{0.1in}
\setlength{\evensidemargin}{0.1in}
\setlength{\textheight}{8.6in}
\setlength{\topmargin}{-0.3in}
\setlength{\headheight}{14pt}
\setlength{\headsep}{22pt}
\setlength{\parindent}{1.2em}
\setlength{\parskip}{3pt}
\newcommand{\R}{\mathbb R}
\newcommand{\N}{\mathbb N}
\newcommand{\Nzero}{\mathbb N_0}
\newcommand{\gra}{\operatorname{gra}}
\newcommand{\dom}{\operatorname{dom}}
\newcommand{\sgn}{\operatorname{sgn}}
\newcommand{\pair}[2]{\langle #1,#2\rangle}
\newtheorem{theorem}{Theorem}
\newtheorem{proposition}{Proposition}
\makeatletter
\def\ps@internal{%
  \def\@oddhead{\small\bfseries DO NOT UPLOAD\hfil FIRST DRAFT / source r1}%
  \let\@evenhead\@oddhead
  \def\@oddfoot{\small Independent review pending\hfil\thepage}%
  \let\@evenfoot\@oddfoot}
\makeatother
\pagestyle{internal}
\author{Weifeng Yang\\
% Yunnan Earthquake Agency, Kunming, China\\
\texttt{ywf841673182@gmail.com}
}
\makeatletter
\newcommand{\preservedAuthor}{\begin{tabular}{c}\@author\end{tabular}}
\makeatother
\begin{document}

\begin{center}
{\Large\bfseries A nonmaximal sum of maximally monotone operators\\[4pt]
with a positive rank-one factor}\\[10pt]
\preservedAuthor\\[10pt]
September 7, 2026
\end{center}

\noindent\textbf{Internal review copy.}
The results below are presented for independent mathematical review.
Human sign-off and public release remain pending. This draft does not claim
a complete machine-checked proof.

\begin{abstract}
We construct two maximally monotone operators on a real Banach space
whose domains satisfy Rockafellar's interior-domain condition and whose
sum is not maximally monotone. The second operator is bounded, linear,
positive, rank one, and defined on the whole space. We first construct
the operators on $c_0$ and then realize the same conclusion on standard
$\ell^1$ with its full dual $\ell^\infty$. The construction uses countably
many copies of a classical triangular operator and constrains their sums
by a Lipschitz curve. Every parameter in this curve has a finitely
supported tail, and its limiting tail belongs to $\ell^2\setminus\ell^1$.
We prove maximality by reducing every point of the monotone polar to a
scalar inequality and excluding the missing parameter interval through
this failure of summability. The first operator omits zero and satisfies
a uniform lower bound on the pairing divided by the primal norm.
The positive rank-one perturbation makes the pairing strictly positive
on every point of the sum graph, so that $(0,0)$ gives a proper monotone
extension. We also compute a positive Fitzpatrick split value at this point.
\end{abstract}

\noindent\textbf{Keywords:} maximally monotone operator; sum theorem;
nonreflexive Banach space; monotone polar; rank-one operator.

\section{Introduction}
\label{sec:introduction}
Let $X$ be a real Banach space, let $X^*$ be its continuous dual, and let
$A,B:X\rightrightarrows X^*$ be maximally monotone operators. Their pointwise
sum is defined by $(A+B)x=\{a+b:a\in Ax,\ b\in Bx\}$. Rockafellar's sum
theorem \cite{rockafellar1970} states that $A+B$ is maximally monotone when
$X$ is reflexive and
\begin{equation}
\dom A\cap\operatorname{int}(\dom B)\ne\varnothing.
\label{eq:original-cq}
\end{equation}
The unrestricted sum question asks whether Eq.~\eqref{eq:original-cq}
alone is sufficient on an arbitrary real Banach space. The interior here
is taken in the norm topology of $X$, and maximality is taken in the full
pair $X\times X^*$.

Several positive results extend the theorem by imposing structure on one
of the factors. Yao \cite{yao2010} proves maximality when the first factor
is of type~(FPV) and the second factor has full domain. Borwein and Yao
\cite{borweinyao2012} prove the sum theorem under
Eq.~\eqref{eq:original-cq} when the first factor is a linear relation.
The factor order in the latter qualification matters: placing a
full-domain linear operator first would require interior points in the
domain of the other factor. These results motivate the question of
whether an elementary full-domain second factor can already produce a
nonmaximal sum with a nonlinear first factor.

Nonreflexive constructions provide a concrete setting for this question.
Bueno and Svaiter \cite{buenosvaiter2011} construct a non-type~(D)
maximally monotone operator on $c_0$. Bauschke, Borwein, Wang, and Yao
\cite{bauschkeetal2011} develop a family of triangular operators and
further nonlinear examples. In particular, their Example~4.1 contains
the single-block triangular map used below. These earlier ingredients
give exact pairing identities. To obtain the sum behavior considered
here, we impose an additional constraint on the sums over countably many
blocks and verify maximality for every point of the full monotone polar.

We define this constraint through a $1$-Lipschitz map
$F:J\to\ell^1(\Nzero)\subset\ell^2(\Nzero)$, where
$J=(-\infty,-1)\cup(1,\infty)$. Its positive-index coordinates have finite
support for each actual parameter. At both deleted endpoints they tend,
coordinatewise and in $\ell^2$, to
$(1/(4n))_{n\ge1}$, which is not in $\ell^1$.
The triangular pairing identity first proves monotonicity. The complete
affine fibres of the graph then reduce the condition for an arbitrary
polar point to
\begin{equation}
\|F(t)-m_p\|_2^2+\nu^2\le(t-\tau)^2
\qquad(t\in J),
\label{eq:polar-intro}
\end{equation}
where $m_p\in\ell^1(\Nzero)$ is determined by its dual coordinate and
$\tau,\nu\in\R$. If $|\tau|>1$, the actual test $t=\tau$ returns a graph
point. If $|\tau|\le1$, finite-coordinate limits at the two endpoints
force the positive-index coordinates of $m_p$ to equal $1/(4n)$.
Their nonsummability excludes this second case and proves maximality.

The resulting first factor $A$ has nonempty domain at positive distance
from zero. Its construction also gives
$\langle x,a\rangle>-\sigma$ on $\gra A$, with
$\sigma=\pi^2/96<1/8$. We choose
$Px=\langle x,g\rangle g$ for a specified nonzero $g\in(c_0)^*$.
The same pairing identity gives
$\langle x,a+Px\rangle>1-\sigma$ on every sum row.
Thus $(0,0)$ is monotonically related to the entire sum and is not in
its graph. The full domain of $P$ makes
Eq.~\eqref{eq:original-cq} immediate from nonemptiness of $\dom A$.

The construction on $c_0$ already addresses the unrestricted question.
We also transfer it to standard $\ell^1$ through an explicitly specified
surjection $Q:\ell^1\to c_0$. The pullback $T=Q^*AQ$ is proved maximally
monotone in $\ell^1\times\ell^\infty$: tests along $\ker Q$ force every
polar dual coordinate to factor through $Q$, and the remaining inequality
is exactly the one for $A$. This argument uses a controlled preimage for
each point, not a bounded linear right inverse.

The principal results are stated in Section~\ref{sec:statements}.
Section~\ref{sec:definitions} defines the graphs and notation.
Sections~\ref{sec:c0} and \ref{sec:l1} give the complete proofs on $c_0$
and $\ell^1$, respectively. Section~\ref{sec:domain-comparison} records
the resulting domain geometry and its relation to a published
almost-convexity statement. Conclusions are given in
Section~\ref{sec:conclusion}. Appendix~\ref{sec:split} computes the
Fitzpatrick split values; these calculations are not needed to prove
nonmaximality of either sum.

% SOURCE: THEOREM_STATEMENT_FROZEN.md:25-42.
\section{Main results}
\label{sec:statements}
Throughout the statements, $\sigma=\pi^2/96<1/8$ and $\mathcal V$ is
the quantity in Eq.~\eqref{eq:D3}. We state the standard $\ell^1$ result
first and then the $c_0$ construction on which its proof is based.
\renewcommand{\thetheorem}{L}
\begin{theorem}[Standard $\ell^1$ realization]\label{thm:L}
On $Z=\ell^1(\N)$ with its usual norm and full dual $Z^*=\ell^\infty(\N)$,
let $T:Z\rightrightarrows Z^*$ and $B:Z\to Z^*$ be defined by Eq.~\eqref{eq:D10}.
Then:
\begin{enumerate}
\item $T$ is maximally monotone in the full pair $Z\times Z^*$.
\item $B$ is nonzero, bounded, positive, rank-one, and maximally monotone,
with $\dom B=Z$. In particular,
$\dom T\cap\operatorname{int}(\dom B)\ne\varnothing$.
\item $T(0)=\varnothing$, and every $(u,b)\in\gra T$ satisfies
$\|u\|_1>1/2$ and $\pair{u}{b}\ge-2\sigma\|u\|_1$.
Thus $\mathcal V(\gra T)\le2\sigma<1/4$.
\item Every $(u,c)\in\gra(T+B)$ satisfies
$\pair{u}{c}>1-\sigma>7/8$.
\item $(0,0)$ belongs to the full monotone polar of $\gra(T+B)$ and does
not belong to $\gra(T+B)$. Hence $T+B$ is not maximally monotone.
\end{enumerate}
\end{theorem}

% SOURCE: THEOREM_STATEMENT_FROZEN.md:6-23.
\renewcommand{\thetheorem}{C}
\begin{theorem}[$c_0$ core]\label{thm:C}
Let $I=\Nzero\times\N$ and $Y=c_0(I)$ with its usual supremum norm and
full dual $Y^*=\ell^1(I)$. Let $A:Y\rightrightarrows Y^*$ and the bounded
positive rank-one map $P:Y\to Y^*$ be defined by Eq.~\eqref{eq:D8}.
Then:
\begin{enumerate}
\item $A$ and $P$ are maximally monotone in $Y\times Y^*$.
\item $\dom P=Y$ and $\dom A\ne\varnothing$, so
$\dom A\cap\operatorname{int}(\dom P)\ne\varnothing$.
\item Every $(x,a)\in\gra A$ satisfies $\|x\|_\infty>1/2$ and
$\pair{x}{a}\ge-2\sigma\|x\|_\infty$, where $\sigma=\pi^2/96<1/8$.
\item Every $(x,b)\in\gra(A+P)$ satisfies
$\pair{x}{b}>1-\sigma>7/8$.
\item $(0,0)$ is monotonically related to the entire graph of $A+P$ and
does not belong to that graph. Hence $A+P$ is not maximally monotone.
\end{enumerate}
\end{theorem}

% SOURCE: THEOREM_STATEMENT_FROZEN.md:20-23,39-49; DEPENDENCY_DAG.md.
\noindent Theorem~\ref{thm:C} supplies a single specified real Banach space
for the unrestricted sum question with the ordered condition
$\dom A\cap\operatorname{int}(\dom P)\ne\varnothing$.
Theorem~\ref{thm:L} uses the same $A$ core, a surjective
$Q:\ell^1\to c_0$, maximality of the pullback in the full dual pair, and
the rank-one factor. Appendix~\ref{sec:split} proves the additional
identities $\Delta_{A,P}(0,0)=\Delta_{T,B}(0,0)=\sigma$, each uniquely
minimized at the zero dual split.

% SOURCE: DEFINITIONS_FROZEN.md:3-22.
\section{Definitions}
\label{sec:definitions}
All scalars are real. Put $\N=\{1,2,\ldots\}$, $\Nzero=\{0,1,\ldots\}$,
and $I=\Nzero\times\N$. The space $Y=c_0(I)$ is global: for each
$\epsilon>0$, only finitely many coordinates have modulus at least
$\epsilon$. It has norm $\|\cdot\|_\infty$ and full dual $\ell^1(I)$,
with $\pair{x}{a}=\sum_{b,j}x_{b,j}a_{b,j}$.
The same pairing is used when the first argument lies in $\ell^\infty(I)$.
The unit coordinate vectors are $e_{b,j}$; the ambient space follows their use.

For a graph $G\subset X\times X^*$, define
\begin{equation}
G^\mu=\{(z,p)\in X\times X^*:
 \pair{z-x}{p-a}\ge0\quad\forall(x,a)\in G\}.
\tag{D1}\label{eq:D1}
\end{equation}
The graph is monotone when every two graph points have a nonnegative bracket,
and maximally monotone when it has no proper monotone extension in this same
full dual pair. For a monotone $G$, $G=G^\mu$ is equivalent to maximality.
Domain and range refer to actual graph points. Every inequality quantified
over the graph includes every value at a multivalued primal.

% SOURCE: DEFINITIONS_FROZEN.md:24-41.
For a maximally monotone $A$, the sign conventions are
\begin{equation}
\begin{aligned}
\Phi_A(d,p)&=F_A(d,p)-\pair{d}{p}
 =\sup_{(y,a)\in\gra A}\pair{d-y}{a-p},\\
\Delta_{A,B}(x,x^*)&=\inf_{u^*\in X^*}
 [\Phi_A(x,x^*-u^*)+\Phi_B(x,u^*)].
\end{aligned}
\tag{D2}\label{eq:D2}
\end{equation}
For the present graphs, which have no zero primal, define
\begin{equation}
\mathcal V(G)=\sup_{(x,a)\in G}
 \frac{\max\{0,-\pair{x}{a}\}}{\|x\|}.
\tag{D3}\label{eq:D3}
\end{equation}

% SOURCE: DEFINITIONS_FROZEN.md:44-59.
We index coordinates by blocks. The map $m$ records the sum in each block,
and $E$ applies a triangular operator within that block. Its single-block
formula is the $\alpha=(1,1,\ldots)$ case of
\cite[Example~4.1]{bauschkeetal2011}; the additional constraint below
couples the block sums through $F$.
For $a\in\ell^1(I)$, define the bounded maps
\begin{equation}
(ma)_b=\sum_{j\ge1}a_{b,j},\qquad
(Ea)_{b,j}=a_{b,j}+2\sum_{k>j}a_{b,k}.
\tag{D4}\label{eq:D4}
\end{equation}
Here $m:\ell^1(I)\to\ell^1(\Nzero)\subset\ell^2(\Nzero)$ and
$E:\ell^1(I)\to Y$. Put
\begin{equation}
\begin{aligned}
g&=e_{0,1}-e_{0,2}\in Y^*,&
h&=Eg=-e_{0,1}-e_{0,2}\in Y,\\
r(a)&=\pair{h}{a}=-a_{0,1}-a_{0,2},&
La&=-Ea+r(a)h.
\end{aligned}
\tag{D5}\label{eq:D5}
\end{equation}

% SOURCE: DEFINITIONS_FROZEN.md:62-78.
For $s>0$ and $n\ge1$, write
\begin{equation}
\omega_s(n)=\frac{\max\{1-sn^{1/4},0\}}{4n},\quad
W(s)=\sum_{n\ge1}\omega_s(n)^2,\quad
\sigma=\sum_{n\ge1}\frac1{16n^2}=\frac{\pi^2}{96}.
\tag{D6}\label{eq:D6}
\end{equation}
Actual $\omega_s$ has finite support. The symbol $\omega_0$ denotes only
the $\ell^2$ comparison vector $(1/(4n))_{n\ge1}$ and is never an actual
graph parameter. Define
\begin{equation}
\begin{aligned}
J&=(-\infty,-1)\cup(1,\infty),\\
F(t)&=(\sgn t,\omega_{|t|-1})\in\ell^1(\Nzero)\quad(t\in J),\\
D&=\{a\in\ell^1(I):r(a)\in J,\ ma=F(r(a))\}.
\end{aligned}
\tag{D7}\label{eq:D7}
\end{equation}

% SOURCE: DEFINITIONS_FROZEN.md:81-95.
The operators and kernel are
\begin{equation}
\begin{aligned}
\gra A&=\{(La,a):a\in D\},\\
K&=\{k\in\ell^1(I):mk=0,\ r(k)=0\},\\
Px&=\pair{x}{g}g\quad(x\in Y).
\end{aligned}
\tag{D8}\label{eq:D8}
\end{equation}
For each $t\in J$, the following is an actual label, not the entire graph:
\begin{equation}
a^t=-t e_{0,1}+(t+\sgn t)e_{0,3}
       +\sum_{n\ge1}\omega_{|t|-1}(n)e_{n,1}.
\tag{D9}\label{eq:D9}
\end{equation}

% SOURCE: DEFINITIONS_FROZEN.md:98-112.
For standard $\ell^1$, let $Z=\ell^1(\N)$ and $Z^*=\ell^\infty(\N)$,
with their usual pairing. Fix once and for all an enumeration $(q_n)$ of
all finitely supported rational vectors in the closed unit ball of $Y$.
Define
\begin{equation}
\begin{aligned}
Qu&=\sum_{n\ge1}u_nq_n,& (Q^*a)_n&=\pair{q_n}{a},\\
\gra T&=\{(u,Q^*a):a\in D,\ Qu=La\},\\
f&=Q^*g,& Bu&=\pair{u}{f}f.
\end{aligned}
\tag{D10}\label{eq:D10}
\end{equation}
The proof establishes surjectivity with a norm-$2$ lift. No bounded linear
right inverse or complemented kernel is assumed. The polar for $T$ and
$T+B$ uses all $\ell^\infty$ labels.

\section{Proof of the $c_0$ core}
\label{sec:c0}
% SOURCE: C0_CORE_PROOF_FROZEN.md:8-35.
\subsection{C1. Types and the bilinear identity}
\label{sec:C1}
For $a\in\ell^1(I)$, $\|ma\|_1\le\|a\|_1$ and
$\|Ea\|_\infty\le2\|a\|_1$. If $a$ has finite support, then $Ea$ has
finite support: only finitely many blocks occur, and within a block no
output occurs after the last nonzero input coordinate. Approximation by
finite-support vectors and the bound prove $Ea\in c_0(I)$ globally.
Thus $L$ also maps $\ell^1(I)$ boundedly into $c_0(I)$.

For $a,c\in\ell^1(I)$, absolute convergence permits exchanging the double
sums. In each block the diagonal terms and the two triangular
off-diagonal sums give
\begin{equation}
\pair{Ea}{c}+\pair{Ec}{a}
 =2\sum_b\Big(\sum_j a_{b,j}\Big)\Big(\sum_j c_{b,j}\Big)
 =2\pair{ma}{mc}_{\ell^2}.
\tag{C1}\label{eq:C1}
\end{equation}
The absolute sum of all products is at most $\|a\|_1\|c\|_1$ before the
factor $2$, so this is an identity for all labels. Since $mg=0$, $h=Eg$,
and $r(g)=0$, Eq.~\eqref{eq:C1} implies $\pair{Ea}{g}=-r(a)$.
Consequently
\begin{equation}
\begin{aligned}
\pair{La}{a}&=r(a)^2-\|ma\|_2^2,&
\pair{La}{g}&=r(a),\\
\pair{La-Lc}{a-c}&=(r(a)-r(c))^2-\|ma-mc\|_2^2.
\end{aligned}
\tag{C2}\label{eq:C2}
\end{equation}

% SOURCE: C0_CORE_PROOF_FROZEN.md:38-64.
\subsection{C2. The actual curve and nonempty fibres}
\label{sec:C2}
For each $s>0$, $\omega_s(n)=0$ whenever $n\ge s^{-4}$. Also
$0\le\omega_s(n)\le1/(4n)$. Hence $W(s)\le\sigma<1/8$, and dominated
convergence in $\ell^2$ gives $\omega_s\to\omega_0$ and $W(s)\to\sigma$
as $s$ decreases to zero. The comparison vector $\omega_0$ is not in
$\ell^1$ by divergence of the harmonic series. No endpoint is added to $J$.

Since $v\mapsto\max\{v,0\}$ is $1$-Lipschitz on $\R$,
\begin{equation}
\|\omega_s-\omega_q\|_2^2\le c|s-q|^2,\qquad
c=\frac1{16}\sum_{n\ge1}n^{-3/2}<\frac3{16}<1.
\tag{C3}\label{eq:C3}
\end{equation}
The strict numerical bound follows from
$\sum_{n\ge1}n^{-3/2}<1+\int_1^\infty x^{-3/2}\,dx=3$.
For $t,u$ on the same branch of $J$, Eq.~\eqref{eq:C3} gives
$\|F(t)-F(u)\|_2\le|t-u|$. For $t=1+s$ and $u=-1-q$, the squared
distance is at most $4+c(s-q)^2\le(2+s+q)^2$, and the opposite
orientation is identical. Thus $F$ is $1$-Lipschitz on all of $J$.

The finite label $a^t$ in Eq.~\eqref{eq:D9} satisfies $r(a^t)=t$ and
$ma^t=F(t)$, proving nonemptiness for every actual $t$. The complete
fibre $\{a:r(a)=t,\ ma=F(t)\}$ is precisely $a^t+K$.
Applying Eq.~\eqref{eq:C2} and the Lipschitz inequality to any two labels
in $D$ proves monotonicity of the entire $A$.

% SOURCE: C0_CORE_PROOF_FROZEN.md:66-105.
\subsection{C3. Whole-polar reduction}
\label{sec:C3}
Take an arbitrary $(x,p)$ in the monotone polar, in the full pair
$Y\times Y^*$. Expansion using Eq.~\eqref{eq:C1} gives, for every
$a\in D$ and $t=r(a)$,
\begin{equation}
\begin{aligned}
\pair{x-La}{p-a}
 &=\pair{x}{p}-\pair{x+Ep}{a}\\
 &\quad+2\pair{ma}{mp}-t r(p)+t^2-\|ma\|_2^2.
\end{aligned}
\tag{C4}\label{eq:C4}
\end{equation}
Replace $a$ by $a+\theta k$ for $k\in K$ and all $\theta\in\R$.
All terms except $-\theta\pair{x+Ep}{k}$ are constant.
Nonnegativity for all $\theta$ forces $\pair{x+Ep}{k}=0$ for every $k\in K$.

We now prove $K^\perp\cap c_0(I)=\R h$ with no dual-space enlargement.
Let $z\in c_0(I)$ globally annihilate $K$. In a block $b\ge1$, every
difference $e_{b,j}-e_{b,k}$ is in $K$. Thus $z$ is constant in that
entire block, and global $c_0$ forces that constant to vanish. In block
zero, differences between indices $j,k\ge3$ belong to $K$, so the common
tail value is also zero. Finally, $g=e_{0,1}-e_{0,2}$ belongs to $K$,
forcing $z_{0,1}=z_{0,2}$. It follows that $z$ is a scalar multiple of
$h$. Conversely, $\pair{h}{k}=r(k)=0$ proves that every multiple of $h$
annihilates $K$. Therefore $x=-Ep+vh$ for some $v\in\R$.

Put $r_p=r(p)$, $m_p=mp$, $\tau=(v+r_p)/2$, and $\nu=(v-r_p)/2$.
Substitution into Eq.~\eqref{eq:C4}, including
$\pair{Ep}{p}=\|mp\|_2^2$, yields the exact bracket
\begin{equation}
\pair{x-La}{p-a}
 =(t-\tau)^2-\nu^2-\|F(t)-m_p\|_2^2.
\tag{C5}\label{eq:C5}
\end{equation}
Every $t\in J$ has an actual $a^t$, and Eq.~\eqref{eq:C5} is constant
over its complete fibre. Thus the full polar condition is exactly
\begin{equation}
\|F(t)-m_p\|_2^2+\nu^2\le(t-\tau)^2
\qquad\forall t\in J.
\tag{C6}\label{eq:C6}
\end{equation}

% SOURCE: C0_CORE_PROOF_FROZEN.md:108-137.
\subsection{C4. Exhaustion of all polar points}
\label{sec:C4}
If $|\tau|>1$, use the actual $t=\tau$ in Eq.~\eqref{eq:C6}. It forces
$\nu=0$ and $m_p=F(\tau)$. Then $v=r_p=\tau$, so $p\in D$ and $x=Lp$,
yielding graph membership.

If $|\tau|\le1$, write $m_p=(b,z_1,z_2,\ldots)$. For each fixed positive
integer $N$, discard the nonnegative tail of Eq.~\eqref{eq:C6}, then
take actual $t\to1$ from above and $t\to-1$ from below.
Finite-coordinate convergence gives
\begin{equation}
\begin{aligned}
(b-1)^2+\nu^2+\sum_{n=1}^N(z_n-1/(4n))^2&\le(1-\tau)^2,\\
(b+1)^2+\nu^2+\sum_{n=1}^N(z_n-1/(4n))^2&\le(1+\tau)^2.
\end{aligned}
\tag{C7}\label{eq:C7}
\end{equation}
Weight these by $(1+\tau)/2$ and $(1-\tau)/2$, respectively. Both
weights are nonnegative, including $\tau=1$ and $\tau=-1$. Expanding
the two squares gives
\begin{equation}
(b-\tau)^2+\nu^2+\sum_{n=1}^N(z_n-1/(4n))^2\le0.
\tag{C8}\label{eq:C8}
\end{equation}
For every $N$, all these nonnegative terms vanish. Thus $z_n=1/(4n)$ for
all $n$, contradicting $m_p\in\ell^1(\Nzero)$. No unattainable endpoint
has been used as a graph row, and no infinite interchange is required in
Eq.~\eqref{eq:C7}. This case is impossible. All polar points belong to
$A$, and monotonicity gives the reverse inclusion. Hence $A$ is maximally
monotone in $Y\times Y^*$.

% SOURCE: C0_CORE_PROOF_FROZEN.md:139-151.
\subsection{C5. Missing zero and the whole-row radial bound}
\label{sec:C5}
For every $a\in D$, $t=r(a)$ satisfies $|t|>1$. By Eq.~\eqref{eq:C2}
and $\|g\|_1=2$,
\begin{equation}
\|La\|_\infty\ge|t|/2>1/2,\qquad
\pair{La}{a}=t^2-1-W(|t|-1)>-\sigma.
\tag{C9}\label{eq:C9}
\end{equation}
It follows that $A(0)=\varnothing$ and
$\pair{La}{a}\ge-2\sigma\|La\|_\infty$ for every actual graph label.
In particular $\mathcal V(\gra A)\le2\sigma<1/4$. No compactness,
bounded dual range, or truncation-to-seed hypothesis enters this deduction.

% SOURCE: C0_CORE_PROOF_FROZEN.md:153-181.
\subsection{C6. The rank-one factor and the sum}
\label{sec:C6}
The map $P$ is everywhere defined, linear, bounded with
$\|P\|\le\|g\|_1^2=4$, nonzero because $Pe_{0,1}=g$, and has range $\R g$.
Its monotonicity follows from $\pair{x-y}{Px-Py}=\pair{x-y}{g}^2$.

For completeness, suppose $(z,z^*)$ is related to its entire graph.
Testing at $z+\lambda v$ for any $v\in Y$ and all $\lambda\in\R$ gives
\begin{equation}
-\lambda\pair{v}{z^*-Pz}+\lambda^2\pair{v}{g}^2\ge0.
\tag{C10}\label{eq:C10}
\end{equation}
If the linear coefficient is nonzero, a sufficiently small $\lambda$ of
the opposite sign makes the expression negative. Hence $z^*=Pz$, proving maximality.

The actual label $a^2=-2e_{0,1}+3e_{0,3}$ gives
$La^2=-6e_{0,1}-8e_{0,2}-3e_{0,3}$. Thus
$\dom A\cap\operatorname{int}(\dom P)\ne\varnothing$, with
$\operatorname{int}(\dom P)=Y$. Every sum point has the form
$(La,a+tg)$, $a\in D$ and $t=r(a)$, and
\begin{equation}
\pair{-La}{-a-tg}=2t^2-1-W(|t|-1)>1-\sigma>7/8.
\tag{C11}\label{eq:C11}
\end{equation}
This proves that $(0,0)$ is related to every sum point. The sum is
monotone as the sum of two monotone operators; it has no zero primal by
Eq.~\eqref{eq:C9}. Adding $(0,0)$ is therefore a proper monotone extension.
This completes the proof of Theorem~\ref{thm:C}.

\section{Proof of the standard $\ell^1$ realization}
\label{sec:l1}
% SOURCE: STANDARD_L1_STRENGTHENING_FROZEN.md:7-38.
\subsection{L1. An actual quotient}
\label{sec:L1}
Finitely supported rational vectors in the closed unit ball of $Y$ are
countable and dense in that ball: truncate a $c_0$ vector and approximate
each retained coordinate by a rational inside $[-1,1]$. Fix an enumeration
of all these vectors as $(q_n)$. For $u\in Z=\ell^1(\N)$, the sum defining
$Qu$ converges absolutely in the Banach space $Y$, and
$\|Qu\|_\infty\le\|u\|_1$. Coordinate unit vectors occur among the $q_n$,
so $\|Q\|=1$.

For an arbitrary $x\in Y$, set $r_0=x$. If $r_k\ne0$, choose an actual
$q_{n_k}$ with $\|q_{n_k}-r_k/\|r_k\|_\infty\|_\infty<1/2$ and set
\begin{equation}
c_k=\|r_k\|_\infty,\qquad r_{k+1}=r_k-c_kq_{n_k}.
\tag{L1}\label{eq:L1}
\end{equation}
If a residual vanishes, stop. Otherwise
$\|r_k\|_\infty\le2^{-k}\|x\|_\infty$ and $\sum_k c_k\le2\|x\|_\infty$.
Define $u_n=\sum_{k:n_k=n}c_k$. Repeated indices are combined.
Positivity of the $c_k$ and summability give $u\in\ell^1$ with
$\|u\|_1=\sum_k c_k\le2\|x\|_\infty$.
Telescoping Eq.~\eqref{eq:L1}, then absolute convergence under regrouping,
proves
\begin{equation}
Qu=x,\qquad \|u\|_1\le2\|x\|_\infty.
\tag{L2}\label{eq:L2}
\end{equation}
For $x=0$, use $u=0$. This proves surjectivity with the asserted lifting
bound. The construction does not assert that $x\mapsto u$ is linear or
continuous, nor that $\ker Q$ is complemented.

% SOURCE: STANDARD_L1_STRENGTHENING_FROZEN.md:40-88.
\subsection{L2. The adjoint and all full-dual polar labels}
\label{sec:L2}
Absolute convergence gives $\pair{u}{Q^*a}=\pair{Qu}{a}$ and
$\|Q^*a\|_\infty\le\|a\|_1$. For each finite subset of $I$, its sign
vector for $a$ is a finitely supported rational vector of norm at most
one and occurs in the enumeration. Taking these finite sign tests proves
\begin{equation}
\|Q^*a\|_\infty=\|a\|_1\qquad(a\in\ell^1(I)).
\tag{L3}\label{eq:L3}
\end{equation}
In particular, $Q^*$ is injective. The graph of $T$ is nonempty by
Eq.~\eqref{eq:L2} and nonemptiness of $D$. It is monotone because the
bracket of any two rows $(u,Q^*a)$ and $(v,Q^*c)$ equals
$\pair{La-Lc}{a-c}$, which is nonnegative by the $c_0$ core.

Let $(v,v^*)$ be a polar point with $v^*$ an arbitrary element of
$\ell^\infty(\N)$. Each actual fibre of $T$ contains $u+\lambda k$ for
every $k\in\ker Q$ and every $\lambda\in\R$. Testing those rows forces
$\pair{k}{v^*}=0$ for every such $k$. Define a functional on $Y$ by
\begin{equation}
\lambda(x)=\pair{u}{v^*}\quad\text{when }Qu=x.
\tag{L4}\label{eq:L4}
\end{equation}
Two preimages differ by $\ker Q$, proving well-definedness. Linear
combinations of preimages prove linearity. For each $x$, choose the
preimage in Eq.~\eqref{eq:L2}; then
$|\lambda(x)|\le2\|v^*\|_\infty\|x\|_\infty$. Thus $\lambda$ belongs to
the continuous dual of $c_0(I)$.

Explicitly put $p_i=\lambda(e_i)$. Finite sign tests give
$\sum_{i\in H}|p_i|\le\|\lambda\|$ for every finite $H$, whence
$p\in\ell^1(I)$. Equality $\lambda(x)=\pair{x}{p}$ holds first for
finite-support $x$ and then for every $c_0$ vector by density and continuity.
Therefore $v^*=Q^*p$: the two functionals agree on every $u$ because
Eq.~\eqref{eq:L4} applies to $Qu$.

Every original $a\in D$ admits a lift $u$ of $La$ by Eq.~\eqref{eq:L2}.
The full $T$-polar inequality now reads
\begin{equation}
\pair{Qv-La}{p-a}\ge0\qquad\forall a\in D.
\tag{L5}\label{eq:L5}
\end{equation}
The original maximality proof gives $p\in D$ and $Qv=Lp$. Hence
$(v,v^*)=(v,Q^*p)$ is an actual $T$ row. This proves maximality in the
full standard pair $\ell^1\times\ell^\infty$. Membership of $v^*$ in
$\operatorname{range}(Q^*)$ was derived from the actual kernel tests.

% SOURCE: STANDARD_L1_STRENGTHENING_FROZEN.md:90-103.
\subsection{L3. Whole-row bounds}
\label{sec:L3}
For every $(u,Q^*a)\in\gra T$, with $t=r(a)$,
\begin{equation}
\begin{aligned}
\pair{u}{Q^*a}&=\pair{La}{a}
 =t^2-1-W(|t|-1)>-\sigma,\\
\|u\|_1&\ge\|Qu\|_\infty=\|La\|_\infty>1/2.
\end{aligned}
\tag{L6}\label{eq:L6}
\end{equation}
Thus $T(0)=\varnothing$ and $\pair{u}{Q^*a}\ge-2\sigma\|u\|_1$ on
the entire graph, including every $\ker Q$ translate. Consequently,
$\mathcal V(\gra T)\le2\sigma<1/4$. No exact equality for this
norm-dependent supremum is asserted.

% SOURCE: STANDARD_L1_STRENGTHENING_FROZEN.md:105-117.
\subsection{L4. Rank-one map and the ordered intersection condition}
\label{sec:L4}
Let $f=Q^*g$. By Eq.~\eqref{eq:L3}, $\|f\|_\infty=\|g\|_1=2$, so $f$
is nonzero. The map $B:u\mapsto\pair{u}{f}f$ is nonzero, linear, bounded
with norm at most $4$, positive, and rank-one. The same whole-line proof
as Eq.~\eqref{eq:C10}, now testing all $v\in\ell^1$ and labels in
$\ell^\infty$, proves $B$ maximally monotone.
Its domain is all of $Z$, so $\operatorname{int}(\dom B)=Z$.

The vector $q_\circ=-(3/4)e_{0,1}-e_{0,2}-(3/8)e_{0,3}$ occurs as
$q_{n_0}$. For $u_\circ=8e_{n_0}$, $Qu_\circ=La^2$.
Thus $(u_\circ,Q^*a^2)$ is an actual $T$ row, and
$\dom T\cap\operatorname{int}(\dom B)\ne\varnothing$ in the stated
factor order.

% SOURCE: STANDARD_L1_STRENGTHENING_FROZEN.md:119-134.
\subsection{L5. Strict polar witness and graph nonmembership}
\label{sec:L5}
For every row $(u,Q^*a)$ of $T$, let $t=r(a)$. Then
$\pair{u}{f}=\pair{Qu}{g}=\pair{La}{g}=t$.
Consequently each sum row has work
\begin{equation}
\pair{u}{Q^*a+Bu}
 =2t^2-1-W(|t|-1)>1-\sigma>7/8.
\tag{L7}\label{eq:L7}
\end{equation}
The sum is monotone by monotonicity of its factors.
Equation~\eqref{eq:L7} shows that $(0,0)$ is related to all its rows.
Since $T(0)=\varnothing$, $(0,0)$ is not a sum row. Adjoining this point
gives a proper monotone extension, completing the proof of
Theorem~\ref{thm:L} with the standard Banach space, full dual, and ordered
intersection condition. No theorem about sums was used to prove either
factor maximal.

\section{Domain geometry and an almost-convexity statement}
\label{sec:domain-comparison}
The same graph also has a nonconvex domain closure. This consequence
helps identify the relationship between the construction and claims
about the domains of arbitrary maximally monotone operators.

\begin{proposition}\label{prop:domain-gap}
For the operator $A$ of Theorem~\ref{thm:C},
$0\in\operatorname{conv}(\dom A)$ and
$\operatorname{dist}(0,\dom A)\ge1/2$.
In particular, $\overline{\dom A}$ is not convex.
\end{proposition}
\begin{proof}
At $t=2$ and $t=-2$, the tails in Eq.~\eqref{eq:D9} vanish, so that
$a^{-2}=-a^2$. Since $L$ is linear, both $La^2$ and $-La^2$ belong
to $\dom A$. Their midpoint is zero. Equation~\eqref{eq:C9} gives
$\|x\|_\infty>1/2$ for every $x\in\dom A$, and hence the asserted
distance bound. A convex closure containing $La^2$ and $-La^2$ would
contain zero, contradicting this bound.
\end{proof}

Eberhard and Wenczel \cite[Theorem~57]{eberhardwenczel2024} state that
the norm closure of the domain of every maximally monotone operator on
a Banach space is convex. Proposition~\ref{prop:domain-gap}, together
with the maximality proof in Section~\ref{sec:C4}, gives an incompatible
conclusion for the displayed operator. The proof on
\cite[pp.~39--40]{eberhardwenczel2024} uses a small-parameter estimate
from Proposition~56, Eq.~(65), along a family in which the operator and
the evaluated point vary. The estimate for a fixed operator and point
has an object-dependent threshold; it does not supply a threshold
uniform over that family. The additional uniform estimate is needed
for that application. Our construction and maximality argument do not
use this published domain-closure statement or its negation.

There is also a version-specific issue with the separable-Asplund sum
statement in Bo\c{t} and Csetnek~\cite[Corollary~11]{botcsetnek2008}.
Its written hypotheses apply to the $c_0(I)$ construction: this space
is separable Asplund and $\dom P=Y$, so the convexified domain difference
is all of $Y$. The inspected proof invokes their Theorem~3, which asserts
bidual conjugate domination for representatives of every maximally
monotone operator when the dual space is separable. The operator in
\cite[Lemma~2.1]{buenosvaiter2011} shows why that assertion cannot be used.
Write $G$ for the Gossez skew map on $\ell^1$, $e=(1,1,\ldots)$, and
$s(a)=\sum_n a_n$. That operator has graph
$\{(-Ga,a):a\in\ell^1,\ s(a)=0\}$. Its Fitzpatrick function is the
indicator of $x+Gb\in\R e$; since $x\in c_0$ and $Gb$ tends to $-s(b)$,
the finite-value points satisfy $x=-Gb-s(b)e$. For its conjugate on
$\ell^1\times\ell^\infty$, direct substitution therefore gives
\begin{equation}
\begin{aligned}
F_S^{c*}(p,\xi)
 &=\sup_{b\in\ell^1}\pair{Gp-s(p)e+\xi}{b},\\
F_S^{c*}(e_1,e-Ge_1)&=0<1=\pair{e-Ge_1}{e_1}.
\end{aligned}
\label{eq:bidual-interface}
\end{equation}
This calculation refutes the domination assertion used in that proof;
failure of the proof alone does not decide the standalone sum statement.
Theorem~\ref{thm:C} is proved independently of it. Although standard
$\ell^1$ is not Asplund, Theorem~\ref{thm:L} still uses the maximality
of $A$ established in Section~\ref{sec:C4}; passing to $\ell^1$ does not
remove that dependency. The distinction concerns the cited 2008 preprint,
not an assertion about every later version or possible repair.

\section{Conclusions}
\label{sec:conclusion}
We constructed maximally monotone operators on $c_0$ and on standard
$\ell^1$ whose sums with specified positive rank-one full-domain
operators are not maximally monotone. In both spaces, the classical
interior-domain condition holds and $(0,0)$ gives a proper monotone
extension of the sum. The maximality proof combines an exact triangular
pairing identity, complete affine fibres, and finite-coordinate endpoint
limits that force a nonsummable dual label in the only remaining case.
The quotient argument preserves the full dual and transfers the same
sum witness to $\ell^1$. The construction also gives a domain at
positive distance from zero whose convex hull contains zero, together
with the explicit Fitzpatrick split values computed below.

\appendix
\section{Optional all-dual split values}
\label{sec:split}
% SOURCE: C0_CORE_PROOF_FROZEN.md:183-205.
\subsection{C7. The $c_0$ split}
\label{sec:C7}
For $\alpha\in\R$, define $H(\alpha)=\Phi_A(0,\alpha g)$.
Actual realization of each $t$, Eq.~\eqref{eq:C2}, and
Eq.~\eqref{eq:C9} yield
\begin{equation}
\begin{aligned}
H(\alpha)&=\sup_{|t|>1}[1+W(|t|-1)-t^2+\alpha t]
 \le1+\sigma+\alpha^2/4,\\
H(\alpha)&\ge\sigma+|\alpha|,\qquad H(0)=\sigma.
\end{aligned}
\tag{C12}\label{eq:C12}
\end{equation}
The lower bound uses both actual endpoint limits. For $p\in Y^*$,
$\Phi_P(0,p)=\sup_x[\pair{x}{p}-\pair{x}{g}^2]$.
If $p\notin\R g$, it does not annihilate $\ker g$: a coordinate outside
the first two, or the sum of those two coordinates, supplies a
finite-support witness. Scaling that witness makes the supremum infinite.
If $p=\beta g$, all scalar values $\pair{x}{g}$ are attainable, giving
$\Phi_P(0,p)=\beta^2/4$.

No $\Phi_A$ value is minus infinity, since its defining graph is nonempty.
Off-span splits therefore contribute $+\infty$. The all-dual infimum is
\[
\Delta_{A,P}(0,0)=\inf_{\beta\in\R}[H(-\beta)+\beta^2/4]=\sigma,
\]
uniquely at $\beta=0$. This refinement is not needed for
Theorem~\ref{thm:C}.

% SOURCE: STANDARD_L1_STRENGTHENING_FROZEN.md:136-151.
\subsection{L6. The standard $\ell^1$ split}
\label{sec:L6}
For $\alpha\in\R$, Eq.~\eqref{eq:L2}, Eq.~\eqref{eq:L6}, and
$\pair{u}{f}=t$ show $\Phi_T(0,\alpha f)=H(\alpha)$, where $H$ is
Eq.~\eqref{eq:C12}. This uses all rows: their expression depends only on
$t$, and each original $a$ admits a lift.

For $w^*\in\ell^\infty$,
$\Phi_B(0,w^*)=\sup_u[\pair{u}{w^*}-\pair{u}{f}^2]$.
If $w^*$ is not a scalar multiple of $f$, some $u\in\ker f$ satisfies
$\pair{u}{w^*}\ne0$. Otherwise, choose $u_0$ with $\pair{u_0}{f}=1$ and
decompose any $u$ as $(u-\pair{u}{f}u_0)+\pair{u}{f}u_0$, forcing
$w^*=\pair{u_0}{w^*}f$, a contradiction.
Scaling that kernel witness gives $\Phi_B(0,w^*)=+\infty$.
For $w^*=\beta f$, all real values of $\pair{u}{f}$ are attainable and
the value is $\beta^2/4$. Therefore the infimum over all $\ell^\infty$
splits is
\[
\Delta_{T,B}(0,0)=\inf_{\beta\in\R}[H(-\beta)+\beta^2/4]=\sigma,
\]
uniquely at $\beta=0$.

\begin{thebibliography}{9}
\bibitem{rockafellar1970}
R.~T. Rockafellar,
On the maximality of sums of nonlinear monotone operators,
\emph{Transactions of the American Mathematical Society}
\textbf{149} (1970), 75--88.

\bibitem{yao2010}
L. Yao,
The sum of a maximal monotone operator of type (FPV) and a maximal
monotone operator with full domain is maximal monotone,
arXiv:1005.2247v2, 2010.

\bibitem{borweinyao2012}
J.~M. Borwein and L. Yao,
Maximality of the sum of a maximally monotone linear relation and a
maximally monotone operator,
arXiv:1212.4266v1, 2012.

\bibitem{buenosvaiter2011}
O. Bueno and B.~F. Svaiter,
A non-type (D) operator in $c_0$,
arXiv:1103.2349v1, 2011.

\bibitem{bauschkeetal2011}
H.~H. Bauschke, J.~M. Borwein, X. Wang, and L. Yao,
Construction of pathological maximally monotone operators on
non-reflexive Banach spaces,
arXiv:1108.1463v1, 2011.

\bibitem{eberhardwenczel2024}
A. Eberhard and R. Wenczel,
Representative functions, variational convergence and almost convexity,
\emph{Set-Valued and Variational Analysis} \textbf{32} (2024), article~4.
doi:10.1007/s11228-024-00708-4.

\bibitem{botcsetnek2008}
R.~I. Bo\c{t} and E.~R. Csetnek,
Weak regularity conditions for maximal monotonicity in separable
Asplund spaces,
Chemnitz Faculty of Mathematics, Preprint~3, April~1, 2008.
\end{thebibliography}

\end{document}
